\textbf{Пример 1.} Загадано число $7$ (программе оно не известно, она видит только $N = 10$).

\begin{itemize}
\item На запрос <<? 5>> интерактор отвечает \texttt{>}, так как $7 > 5$;
\item на запрос <<? 8>> интерактор отвечает \texttt{<}, так как $7 < 8$;
\item на запрос <<? 6>> интерактор отвечает \texttt{>}, так как $7 > 6$;
\item на запрос <<? 7>> интерактор отвечает \texttt{=}, и программа завершается. Сделано $4$ запроса из разрешённых $10$.
\end{itemize}

\textbf{Пример 2.} При $N = 1$ загадано число $1$. Программа всё равно должна сделать запрос <<? 1>> и получить ответ \texttt{=}.
